% Options for packages loaded elsewhere
\PassOptionsToPackage{unicode}{hyperref}
\PassOptionsToPackage{hyphens}{url}
%
\documentclass[
  a4paper,
]{article}
\usepackage{amsmath,amssymb}
\usepackage{iftex}
\ifPDFTeX
  \usepackage[T1]{fontenc}
  \usepackage[utf8]{inputenc}
  \usepackage{textcomp} % provide euro and other symbols
\else % if luatex or xetex
  \usepackage{unicode-math} % this also loads fontspec
  \defaultfontfeatures{Scale=MatchLowercase}
  \defaultfontfeatures[\rmfamily]{Ligatures=TeX,Scale=1}
\fi
\usepackage{lmodern}
\ifPDFTeX\else
  % xetex/luatex font selection
\fi
% Use upquote if available, for straight quotes in verbatim environments
\IfFileExists{upquote.sty}{\usepackage{upquote}}{}
\IfFileExists{microtype.sty}{% use microtype if available
  \usepackage[]{microtype}
  \UseMicrotypeSet[protrusion]{basicmath} % disable protrusion for tt fonts
}{}
\makeatletter
\@ifundefined{KOMAClassName}{% if non-KOMA class
  \IfFileExists{parskip.sty}{%
    \usepackage{parskip}
  }{% else
    \setlength{\parindent}{0pt}
    \setlength{\parskip}{6pt plus 2pt minus 1pt}}
}{% if KOMA class
  \KOMAoptions{parskip=half}}
\makeatother
\usepackage{xcolor}
\usepackage[margin=22mm]{geometry}
\usepackage{longtable,booktabs,array}
\usepackage{calc} % for calculating minipage widths
% Correct order of tables after \paragraph or \subparagraph
\usepackage{etoolbox}
\makeatletter
\patchcmd\longtable{\par}{\if@noskipsec\mbox{}\fi\par}{}{}
\makeatother
% Allow footnotes in longtable head/foot
\IfFileExists{footnotehyper.sty}{\usepackage{footnotehyper}}{\usepackage{footnote}}
\makesavenoteenv{longtable}
\setlength{\emergencystretch}{3em} % prevent overfull lines
\providecommand{\tightlist}{%
  \setlength{\itemsep}{0pt}\setlength{\parskip}{0pt}}
\setcounter{secnumdepth}{-\maxdimen} % remove section numbering
\usepackage{kotex}
\ifLuaTeX
  \usepackage{selnolig}  % disable illegal ligatures
\fi
\IfFileExists{bookmark.sty}{\usepackage{bookmark}}{\usepackage{hyperref}}
\IfFileExists{xurl.sty}{\usepackage{xurl}}{} % add URL line breaks if available
\urlstyle{same}
\hypersetup{
  hidelinks,
  pdfcreator={LaTeX via pandoc}}

\author{}
\date{}

\begin{document}

\hypertarget{sequential-uxad11uxd559-uxacc4uxc0b0-uxc5d4uxc9c4uxc758-uxd575uxc2ec-uxc6d0uxb9acuxc640-uxc218uxc2dd}{%
\section{Sequential 광학 계산 엔진의 핵심 원리와
수식}\label{sequential-uxad11uxd559-uxacc4uxc0b0-uxc5d4uxc9c4uxc758-uxd575uxc2ec-uxc6d0uxb9acuxc640-uxc218uxc2dd}}

작성일: 2026-10-18 · 대상: \texttt{src/sequential\_optics/}의 독립
Sequential 계산 모델

이 문서는 \textbf{사용자 처방이 광선, 파면, 광량을 거쳐 시뮬레이션
결과가 되는 과정}을 설명한다. 실제 구현의 좌표·단위·정규화에 맞춘 원리
문서다. 세부 오류 사례와 수정 이력은 별도 검증 문서에 둔다. 수식은 UTF-8
Markdown 안에 LaTeX 원문으로 보존했다.

계산의 중심은 실제 비구면 교점과 벡터 Snell 법칙이다. 회절은 그 광선에서
얻은 위상과 명시적인 진폭 모델을 연결한 \textbf{scalar FFT 모델}이다.
EFL/BFL의 근축 정의, chief 주변의 미분적 field curvature, scalar 회절
근사는 각각 다른 목적을 갖는다. 이 구성은 공개 광학 이론에 기반하며,
Zemax의 모든 내부 규약이나 수치 결과와의 일치를 뜻하지 않는다.

\hypertarget{uxc804uxccb4-uxad6cuxc870-uxd558uxb098uxc758-uxcc98uxbc29uxc5d0uxc11c-uxc138-uxc885uxb958uxc758-uxc815uxbcf4uxb97c-uxb9ccuxb4e0uxb2e4}{%
\subsection{1. 전체 구조: 하나의 처방에서 세 종류의 정보를
만든다}\label{uxc804uxccb4-uxad6cuxc870-uxd558uxb098uxc758-uxcc98uxbc29uxc5d0uxc11c-uxc138-uxc885uxb958uxc758-uxc815uxbcf4uxb97c-uxb9ccuxb4e0uxb2e4}}

처방으로부터 다음 세 가지를 연결해 계산한다.

\begin{enumerate}
\def\labelenumi{\arabic{enumi}.}
\tightlist
\item
  \textbf{광선의 위치와 방향:} spot, chief mapping, distortion, grid,
  기하 수차.
\item
  \textbf{광선의 위상과 경계장:} OPD, diffraction PSF, complex OTF, MTF.
\item
  \textbf{광선이 운반하는 전력:} 투과율, RI, 원거리 광도, 실제 광원
  영상.
\end{enumerate}

광선의 적분 가중치, 광선 패킷의 전력, 전기장 진폭은 서로 다른 값이다.
기하 spot에서 사용한 pupil 가중치를 그대로 LED 방출 전력으로 해석하거나,
전력의 제곱근을 면적 변환 없이 FFT 진폭으로 사용하지 않는다.

처리 순서는 \textbf{입력 정규화 → 재료·형상 상태 생성 → STOP 설정 →
real-ray 추적 → 기하 분석/회절 분석/광량 분석 → 영상 합성 또는 온도별
반복}이다. 각 분석은 실제로 필요한 데이터만 읽는다.

\hypertarget{uxb2e8uxc704uxc640-uxc88cuxd45c}{%
\subsubsection{1.1 단위와 좌표}\label{uxb2e8uxc704uxc640-uxc88cuxd45c}}

\begin{longtable}[]{@{}
  >{\raggedright\arraybackslash}p{(\columnwidth - 2\tabcolsep) * \real{0.5000}}
  >{\raggedright\arraybackslash}p{(\columnwidth - 2\tabcolsep) * \real{0.5000}}@{}}
\toprule\noalign{}
\begin{minipage}[b]{\linewidth}\raggedright
기호·입력
\end{minipage} & \begin{minipage}[b]{\linewidth}\raggedright
내부 의미
\end{minipage} \\
\midrule\noalign{}
\endhead
\bottomrule\noalign{}
\endlastfoot
\(x,y,z,R,t,S,W\) & mm. \(S\)는 optical path/eikonal, \(W\)는 OPD \\
\(\lambda_0\) & 진공 파장. 입력 nm를 \(10^{-6}\)배하여 mm로 변환 \\
\(n\) & 절대 \textbf{위상} 굴절률; 매질 내 파장은
\(\lambda_m=\lambda_0/n\) \\
\(\mathbf d\) & 단위 진행 방향 \\
field angle & rad; 유한 물체의 field height는 mm \\
focal spatial frequency & cycles/mm = lp/mm \\
angular spatial frequency & cycles/rad; focal 주파수와 별도 설정 \\
\(A_{2m}\) & \(\mathrm{mm}^{1-2m}\) \\
\(\Phi,E,I,L_s\) & 복사 전력, 복사조도, 복사강도, radiance. 절대
입력이면 각각 W, W/mm², W/sr, W/(mm²·sr) \\
\end{longtable}

여기서 광량은 radiometric 양으로 계산한다. W/mm²와 W/sr를 lux와 cd로
해석하지 않는다. Photometric 값에는 파장별 전력과 시감도 가중이 추가로
필요하다.

광축은 \(+z\)이고 각 면의 vertex 위치는 \(z_j\)다. 입력 행의 thickness를
누적해 \(z_{j+1}=z_j+t_j\)로 바꾼다. 렌즈 앞면의 thickness는 중심 두께,
뒷면의 thickness는 다음 면까지의 vertex 간격이다. Active diameter의
절반이 clear radius이며 mechanical radius와 구별한다.

G를 기본 파장으로 선택하되 실제 수치는 \texttt{EngineSettings}에서
관리한다. 배포 예제의 G=550 nm는 변경 가능한 기본값이다. G 단색의
nominal spot/MTF에는 R/B 굴절률, CTE, \(dn/dT\), 흡수율을 요구하지
않는다. 온도나 다른 파장을 요청했을 때만 해당 물성을 사용한다.

\hypertarget{uxc7acuxb8ccuxb97c-uxd568uxc218uxb85c-uxd45cuxd604}{%
\subsubsection{1.2 재료를 함수로
표현}\label{uxc7acuxb8ccuxb97c-uxd568uxc218uxb85c-uxd45cuxd604}}

엔진에 필요한 기본 물성은 \(n(\lambda_0,T)\)이며, 광량에는 전력 흡수계수
\(\mu(\lambda_0,T)\), 열 형상에는 치수 배율 \(s(T;T_0)\)가 추가된다.
굴절률은 공급된 표의 보간 또는 명시된 분산식으로 평가한다. 지원하는
Sellmeier 형식은 다음과 같다. 식 평가 시 파장 단위는 제공 계수와
일치시킨다. {[}R10{]}

\[
n^2(\lambda)-1=\sum_j\frac{B_j\lambda^2}{\lambda^2-C_j}.
\]

Abbe 수 \(V_d=(n_d-1)/(n_F-n_C)\)는 분산의 요약량이다. 이것만으로 임의
RGB 굴절률을 유일하게 복원하지 않는다. 재료의 이름도 수치 물성의
대체물이 아니다.

공기 기준 데이터는 같은 물리 파장과 기준 상태에서
\(n_{\rm abs}=n_{\rm rel}n_{\rm air}\)로 변환한다. 파장도 공기 중
파장으로 주어졌다면 \(\lambda_0=n_{\rm air}\lambda_{\rm air}\)다. 두
변환의 필요 여부를 각각 지정한다. 엔진의 Snell과 OPL에는 동일한 absolute
convention을 사용한다. {[}R11{]}

\hypertarget{uxd615uxc0c1uxc5d0uxc11c-uxc2e4uxc81c-uxad11uxc120uxc744-uxad6cuxd55cuxb2e4}{%
\subsection{2. 형상에서 실제 광선을
구한다}\label{uxd615uxc0c1uxc5d0uxc11c-uxc2e4uxc81c-uxad11uxc120uxc744-uxad6cuxd55cuxb2e4}}

\hypertarget{standarduxc640-even-asphere}{%
\subsubsection{2.1 Standard와 Even
Asphere}\label{standarduxc640-even-asphere}}

\(r^2=x^2+y^2\), curvature \(c=1/R\), conic constant \(k\)일 때 local
sag는 다음과 같다. Standard는 모든 \(A_{2m}=0\)인 경우다. {[}R1{]}

\[
z_s(x,y)=\frac{cr^2}{1+\sqrt{1-(1+k)c^2r^2}}
 +\sum_{m=1}^{6}A_{2m}r^{2m}.
\tag{1}
\]

전역 표면은 \(F(x,y,z)=z-z_j-z_s(x,y)=0\)이다. 법선은 sag 미분으로
구한다.

\[
\mathbf N=\frac{(-\partial_xz_s,-\partial_yz_s,1)}
 {\sqrt{1+(\partial_xz_s)^2+(\partial_yz_s)^2}}.
\tag{2}
\]

구현은 \(q=r^2\)의 다항식으로 sag와 미분을 평가한다. 특히 \(A_2\)도
vertex curvature에 기여한다.

\[
b(q)=\frac{c}{\sqrt{1-(1+k)c^2q}}
 +\sum_{m=1}^{6}2mA_{2m}q^{m-1},\qquad
\nabla z_s=b(q)(x,y),\qquad c_v=c+2A_2.
\tag{3}
\]

\hypertarget{uxc2e4uxc81c-uxad50uxc810uxacfc-uxad74uxc808}{%
\subsubsection{2.2 실제 교점과
굴절}\label{uxc2e4uxc81c-uxad50uxc810uxacfc-uxad74uxc808}}

입사 광선을 \(\mathbf r(t)=\mathbf o+t\mathbf d\), \(\|\mathbf d\|=1\)로
쓴다. 다음 면의 교점은 다음 일변수 방정식의 순방향 첫 유효근이다.

\[
g(t)=o_z+td_z-z_j-z_s(o_x+td_x,o_y+td_y)=0,
\qquad t\ge0.
\tag{4}
\]

해당 면의 실제 sag로 교점을 구한 뒤 clear aperture와 STOP을 적용한다.
기하학적 구간 제한·근 분리·수치 근 찾기는 식 (4)를 푸는 방법이며, 표면을
근축 구면이나 얇은 렌즈로 바꾸는 과정이 아니다.

법선 \(\mathbf N\)을 입사 방향과 반대로 향하게 정렬하고,
\(c_i=-\mathbf d\cdot\mathbf N\ge0\), \(\eta=n_1/n_2\)로 두면 벡터 Snell
식은 다음과 같다. {[}R2{]}

\[
c_t=\sqrt{1-\eta^2(1-c_i^2)},\qquad
\mathbf d'=\eta\mathbf d+(\eta c_i-c_t)\mathbf N.
\tag{5}
\]

균질 구간에서 광선은 직선이고 구간 길이 \(\ell_j\)의 optical path 기여는
\(n_j\ell_j\)다. 굴절 후 위치·방향·누적 OPL을 저장한다.

\[
S=S_{\rm in}+\sum_jn_j\ell_j.
\tag{6}
\]

\hypertarget{stop-aiminguxacfc-uxacf5uxd1b5-uxc785uxc0ac-uxc704uxc0c1}{%
\subsubsection{2.3 STOP aiming과 공통 입사
위상}\label{stop-aiminguxacfc-uxacf5uxd1b5-uxc785uxc0ac-uxc704uxc0c1}}

정규화 pupil 좌표 \(\boldsymbol\rho=(\rho_x,\rho_y)\)의 목표점은 STOP
위의 \((a_{\rm stop}\rho_x,a_{\rm stop}\rho_y)\)다. 발사 위치 또는 방향
\(\boldsymbol\xi\)를 조절하여 실제 STOP 교점이 이 목표에 도달하도록
잔차를 푼다. 초기값에는 근축 행렬을 쓸 수 있으나 최종 잔차는 실제
광선으로 계산한다. {[}R3{]}

\[
\mathbf r_{\rm aim}(\boldsymbol\xi)
 =\mathbf x_{\rm stop}(\boldsymbol\xi)-a_{\rm stop}\boldsymbol\rho=\mathbf0.
\tag{7}
\]

Chief ray는 STOP 중심, marginal ray는 STOP rim을 목표로 한다. EPD 또는
명목 F-number는 기준 상태에서 물리 STOP 반경으로 변환한다. 명목
\(F/\#=|f_{\rm report}|/\mathrm{EPD}\)이며 유한 공액의 working
F-number와 구별한다.

위상은 모든 광선이 같은 입사 파면에서 출발하도록 정의한다. 유한 물점
\(\mathbf o_s\)에서 직접 출발하면 \(S_{\rm in}=0\)이고 첫 구간 길이를
실제로 더한다. 무한 물체 평면파의 고정 발사면에서는

\[
S_{\rm in}(\mathbf a)=n_0\mathbf d_0\cdot(\mathbf a-\mathbf a_{\rm ref}).
\tag{8}
\]

이 입사 위상 항이 서로 다른 발사점의 위상을 연결한다. 발사면이
등위상면이 아닐 때 모든 광선의 초기 OPL을 같게 두면 비축 OPD가 달라진다.

기본 bundle은 STOP 면적의 Gauss 적분을 \(\rho^2\) 방향, 주기적 표본을
방위각 방향에 사용한다. 정규화 가중치 합은 1이다. Chief/rim을 추가 기하
표본으로 저장하되 적분 가중치는 0으로 두어, 통계 적분과 pupil 경계
표현을 함께 유지한다.

\hypertarget{uxad11uxc120uxc73cuxb85cuxbd80uxd130-uxae30uxd558uxd559uxc801-uxc131uxb2a5uxc744-uxad6cuxd55cuxb2e4}{%
\subsection{3. 광선으로부터 기하학적 성능을
구한다}\label{uxad11uxc120uxc73cuxb85cuxbd80uxd130-uxae30uxd558uxd559uxc801-uxc131uxb2a5uxc744-uxad6cuxd55cuxb2e4}}

\hypertarget{spotuxacfc-uxbc29uxd5a5-spot}{%
\subsubsection{3.1 Spot과 방향
spot}\label{spotuxacfc-uxbc29uxd5a5-spot}}

마지막 광선의 위치·방향을 \((\mathbf p_i,\mathbf d_i)\)라 하면 지정된
global-\(z\) 영상면의 교점은 다음과 같다.

\[
\mathbf x_i=\left[\mathbf p_i+
 \frac{z_{\rm IMA}-p_{i,z}}{d_{i,z}}\mathbf d_i\right]_{xy}.
\tag{9}
\]

통과 광선의 pupil 적분 가중치 \(w_i\)를 정규화한
\(\widehat w_i=w_i/\sum_iw_i\)로 centroid와 RMS를 계산한다. 기준점
\(\mathbf x_{\rm ref}\)는 centroid 또는 chief 교점이다.

\[
\overline{\mathbf x}=\sum_i\widehat w_i\mathbf x_i,\qquad
\sigma_{\rm RMS}=\sqrt{\sum_i\widehat w_i
 \|\mathbf x_i-\mathbf x_{\rm ref}\|^2}.
\tag{10}
\]

최대 기하 반경은 성공적으로 추적된 \textbf{모든 기하 표본}에서
\(\max_i\|\mathbf x_i-\mathbf x_{\rm ref}\|\)를 구한다. 이 값은 표본
최대값이고 연속 pupil의 전역 극값 보증은 아니다. 무한대 출력의 spot은
위치 대신 chief-local 기저 \((\mathbf e_1,\mathbf e_2,\mathbf e_3)\)의
방향각을 사용한다.

\[
\alpha_i=\operatorname{atan2}(\mathbf d_i\cdot\mathbf e_1,
 \mathbf d_i\cdot\mathbf e_3),\qquad
\beta_i=\operatorname{atan2}(\mathbf d_i\cdot\mathbf e_2,
 \mathbf d_i\cdot\mathbf e_3).
\tag{11}
\]

\hypertarget{eflbfluxc740-1uxcc28-uxad11uxd559uxb7c9uxc774uxb2e4}{%
\subsubsection{3.2 EFL/BFL은 1차
광학량이다}\label{eflbfluxc740-1uxcc28-uxad11uxd559uxb7c9uxc774uxb2e4}}

이 엔진은 reduced ray \((y,n\theta)\)를 사용한다. 거리 \(t\)의 전파와
vertex 굴절은 다음 행렬이다. {[}R4{]}

\[
T(t,n)=\begin{pmatrix}1&t/n\\0&1\end{pmatrix},\qquad
R(n_1,n_2,c_v)=\begin{pmatrix}1&0\\-(n_2-n_1)c_v&1\end{pmatrix}.
\tag{12}
\]

첫 면 직전에서 마지막 면 직후까지 순서대로 곱한
\(M=\begin{pmatrix}A&B\\C&D\end{pmatrix}\)에 대해 출력 매질 굴절률이
\(n_i\)이면

\[
f'=-\frac{n_i}{C},\qquad
\mathrm{BFL}=-\frac{n_iA}{C},\qquad
f_{\rm report}=-\frac{n_{\rm ref}}{C}.
\tag{13}
\]

BFL의 원점은 마지막 물리면 vertex다. \(n_{\rm ref}\)는 설정한 보고 기준
매질이다. 유한 물체에서는 OBJ→마지막 면의 행렬을 \(M_o\)로 만든 뒤
영상까지의 추가 거리 \(d\)를 포함한 \(T(d,n_i)M_o\)의 \(B=0\) 조건으로
근축 공액을 구한다. 따라서 \(d=-n_iB_o/D_o\), 배율은 \(A_o+dC_o/n_i\)다.
이 정의로 EFL/BFL을 구하는 것이 실제 spot 광선을 근축으로 계산한다는
뜻은 아니다.

\hypertarget{distortion-griduxc640-uxc218uxcc28}{%
\subsubsection{3.3 Distortion, grid와
수차}\label{distortion-griduxc640-uxc218uxcc28}}

물체 좌표 또는 field chart \(\mathbf q\)를 실제 chief-ray 영상 위치로
보내는 함수를 \(\mathbf g(\mathbf q)\)라 한다. 무한대 출력이면
\(\mathbf g\)는 global 방향각이다. 기준 field \(\mathbf q_0\) 주변 실제
chief trace의 미분 \(J_0\)로 기준 mapping을 만든다. {[}R5{]}

\[
\mathbf g_{\rm ref}(\mathbf q)=\mathbf g(\mathbf q_0)+J_0(\mathbf q-\mathbf q_0),
\qquad\Delta\mathbf g=\mathbf g-\mathbf g_{\rm ref}.
\tag{14}
\]

무한 물체의 rectilinear chart는 \((\tan\theta_x,\tan\theta_y)\)이고, 현재 구현의
\texttt{f\_theta}는 성분별 각도 chart \((\theta_x,\theta_y)\)를 사용한다.
이 성분별 chart는 실제 광축각을 사용하는 방사 방향 F-θ와 일반적인 2D field에서 같지 않다.
두 정의를 구별하며, 현재 코드를 방사 방향 F-θ로 해석하지 않는다.
유한 물체의 chart는 object height다. 한
meridian의 distortion은 \(100(h-h_{\rm ref})/h_{\rm ref}\)다. Grid는 각
물체 격자점의 chief mapping을 연결하고, 같은 영상면에 해당 점의
spot/PSF를 배치한다. 온도별 배율 변화도 보려면 기준 온도에 대한 raw
mapping 차이를 보존한다.

기하 T 방향은 field의 방사 방향, S는 그 수직 방향이다. 축상에서는 T=Y,
S=X로 정한다. Ray fan은 각 STOP 직경을 따라 pupil 좌표를 바꾸어 chief에
대한 영상 교점 차이를 표시한다. 축상 구면수차의 longitudinal 표현은

\[
z_{\rm cross}(\rho)=p_z-p_y\frac{d_z}{d_y},\qquad
\mathrm{LSA}(\rho)=z_{\rm cross}(\rho)-z_{\rm ref}.
\tag{15}
\]

\(\rho=0\)의 축 교점은 실제 근축 공액의 극한으로 정의한다. 보고 기준면
\(z_{\rm ref}\)를 이동해도 그 물리적 극한은 바뀌지 않는다. Coma 진단에는
fan \(\Delta(\rho)\)의 짝 성분 \([\Delta(\rho)+\Delta(-\rho)]/2\)를
사용한다. 이는 고차 영향을 포함한 real-ray descriptor이며 Seidel coma
계수와 동일한 정의가 아니다.

Field curvature는 chief 주변 실제 광선을 미분하여 구한다. 기준면
\(z_0\)의 위치와 slope를 \(\mathbf X(\boldsymbol\rho)\),
\(\mathbf b(\boldsymbol\rho)=\mathbf d_{xy}/d_z\)로 두면

\[
J_X=\left.\frac{\partial\mathbf X}{\partial\boldsymbol\rho}\right|_0,
\quad J_b=\left.\frac{\partial\mathbf b}{\partial\boldsymbol\rho}\right|_0,
\quad z_a=z_0-\frac{(J_X)_{aa}}{(J_b)_{aa}},
\quad FC_a=z_a-z_{\rm IMA}\quad(a=T,S).
\tag{16}
\]

마지막 비는 T/S 기저의 해당 대각성분에 적용한다. 이는 chief 주변의
parabasal 초점이며 full-pupil RMS 최소 초점과 다르다. \(z_T-z_S\)는 이
정의의 astigmatic focus separation이다. {[}R5{]}

파면 수차는 뒤에서 정의하는 OPD \(W\)를 Zernike basis에 가중
최소제곱으로 분해한다.

\[
W_i\simeq\sum_j a_jZ_j(\rho_i,\varphi_i),\qquad
\underset{\{a_j\}}{\operatorname{minimize}}
 \sum_iw_i\left(W_i-\sum_ja_jZ_j\right)^2.
\tag{17}
\]

기본 chart는 정규화한 물리 STOP이고, basis는 전체 단위원판에서 RMS=1로
정규화한다. \(m>0\)에서
\(Z_{n,m,c/s}=\sqrt{2(n+1)}R_n^m(\rho)\{\cos m\varphi,\sin m\varphi\}\),
\(m=0\)에서 \(Z_{n,0}=\sqrt{n+1}R_n^0\)다. 구현은 안정적인 Jacobi
다항식으로 같은 \(R_n^m\)을 평가한다. 수차 분해는 진단이며, 이를 이유로
실제 MTF 입력에서 defocus나 tilt를 자동 제거하지 않는다. {[}R6{]}

\hypertarget{opluxc744-uxd68cuxc808-uxacc4uxc0b0uxc6a9-uxd30cuxba74uxc73cuxb85c-uxc5f0uxacb0uxd55cuxb2e4}{%
\subsection{4. OPL을 회절 계산용 파면으로
연결한다}\label{opluxc744-uxd68cuxc808-uxacc4uxc0b0uxc6a9-uxd30cuxba74uxc73cuxb85c-uxc5f0uxacb0uxd55cuxb2e4}}

\hypertarget{uxc2e4uxc81c-uxcd9cuxb825-uxacbduxacc4uxba74uxacfc-eikonal}{%
\subsubsection{4.1 실제 출력 경계면과
eikonal}\label{uxc2e4uxc81c-uxcd9cuxb825-uxacbduxacc4uxba74uxacfc-eikonal}}

모든 렌즈를 통과한 광선들을 마지막 면 이후의 공통 균질 매질 내 평면
\(\Pi\)로 옮긴다. 평면의 법선 \(\mathbf e_3\)는 chief 진행 방향이다.
\(\mathbf e_1\)은 global X를 chief-normal 평면에 투영한 방향,
\(\mathbf e_2=\mathbf e_3\times\mathbf e_1\)로 잡는다.

\[
\mathbf r(u,v)=\mathbf r_0+u\mathbf e_1+v\mathbf e_2,\qquad
t_\Pi=\frac{(\mathbf r_0-\mathbf p)\cdot\mathbf e_3}
 {\mathbf d\cdot\mathbf e_3},\qquad
S_\Pi=S_{\rm last}+n_it_\Pi.
\tag{18}
\]

이 경계면은 실제 출사동공의 근축 위치를 그대로 사용하는 면이 아니다.
실제 전파가 끝난 균질 영역에 두는 계산용 물리 평면이다. 이 위의
\(S_\Pi\)와 전기장 진폭 \(a_\Pi\)가 회절 계산의 입력이다. 기본 경계장은
실제 광선 eikonal을 연결한 단일 scalar 장으로 표현한다. 시간 의존성을
\(e^{-i\omega t}\)로 택하고 진행 위상을 \(e^{+i2\pi S/\lambda_0}\)로
쓴다.

\hypertarget{uxc774uxc0c1-uxc218uxb834-uxad6cuxba74uxd30cuxc640uxc758-uxcc28uxc774uxac00-opduxb2e4}{%
\subsubsection{4.2 이상 수렴 구면파와의 차이가
OPD다}\label{uxc774uxc0c1-uxc218uxb834-uxad6cuxba74uxd30cuxc640uxc758-uxcc28uxc774uxac00-opduxb2e4}}

Chief가 사용자 IMA와 만나는 점을 \(\mathbf q_0\)로 하고
\(\mathbf q_0=\mathbf r_0+L\mathbf e_3\)로 쓴다. \(S_c\)는 경계면에서
chief의 eikonal이다. 이상적으로 \(\mathbf q_0\)에 수렴하는 파와 비교하면

\[
R=\sqrt{L^2+u^2+v^2},\qquad
W(u,v)=S_\Pi(u,v)-S_c+n_i(R-L).
\tag{19}
\]

이상 수렴파는 \(S_\Pi=S_c-n_i(R-L)\)이므로 \(W=0\)이다. 실제 계산은
\(R-L=(u^2+v^2)/(R+L)\)로 평가한다. 위상은 진공 파장 기준
\(\phi=2\pi W/\lambda_0\)다. \(W\)에는 이미 굴절률이 포함되어 있으므로
위상 변환에서 \(n_i\)를 한 번 더 곱하지 않는다.

입사 위상, 실제 OPL, 공통 경계면, reference sphere가 모두 갖추어져야
OPD가 된다. 서로 다른 영상 교점까지의 raw OPL 차이만으로 대체하지
않는다. 기본 계산에서 제거되는 것은 상수 위상인 piston뿐이다.

\hypertarget{scalar-fftuxb85c-psf-otf-mtfuxb97c-uxacc4uxc0b0uxd55cuxb2e4}{%
\subsection{5. Scalar FFT로 PSF, OTF, MTF를
계산한다}\label{scalar-fftuxb85c-psf-otf-mtfuxb97c-uxacc4uxc0b0uxd55cuxb2e4}}

\hypertarget{uxc774-uxad6cuxd604uxc758-fourier-pupil}{%
\subsubsection{5.1 이 구현의 Fourier
pupil}\label{uxc774-uxad6cuxd604uxc758-fourier-pupil}}

일반적인 설명은 복소 pupil \(P=a\exp(i2\pi W/\lambda_0)\)를 Fourier
변환한다고 요약한다. 구현에서는 \textbf{어떤 좌표의 진폭과 어떤 면적
측도인지}까지 지정해야 한다. 본 엔진은 scalar Rayleigh--Sommerfeld
전파를 초점 주변에서 전개하고, 다음 reference-direction-cosine 좌표를
사용한다. 이는 광선 굴절의 근축 근사와 별개인 회절 전파 근사다. {[}R7{]}

\[
\chi=\frac{u}{R},\qquad\eta=\frac{v}{R},\qquad
\gamma=\sqrt{1-\chi^2-\eta^2}=\frac{L}{R},\qquad
du\,dv=\frac{L^2}{\gamma^4}\,d\chi\,d\eta.
\tag{20}
\]

여기서 \((\chi,\eta)\)는 \textbf{reference point를 향하는 기하로 정의한
좌표}다. 수차가 있는 실제 ray의 출사 방향 코사인과 항상 같은 것은
아니다. Boundary 진폭을 \(a_\Pi\)라 할 때 obliquity와 Jacobian을 합친
Fourier 진폭 \(B\)는

\[
B(\chi,\eta)=a_\Pi(u,v)\frac{L}{\gamma^2},\qquad
P(\chi,\eta)=B(\chi,\eta)
 \exp\left(i\frac{2\pi}{\lambda_0}W\right).
\tag{21}
\]

식 (20)--(21)은 본 구현의 좌표변환 유도다. 공개 RS 이론을 사용하는
방식이며 Zemax 내부 mapping을 복제했다는 의미는 아니다. Lens-only 기본
분석은 최초 기준면에서 \textbf{통과 Fourier pupil의 \(B=1\)}을 정의한다.
실제 LED 조명 계산은 7절의 전력 수송으로 \(a_\Pi\)를 얻는다. 두 분석의
진폭은 목적이 다르다.

\hypertarget{uxcd08uxc810-uxc8fcuxbcc0-scalar-uxc804uxd30cuxc640-psf}{%
\subsubsection{5.2 초점 주변 scalar 전파와
PSF}\label{uxcd08uxc810-uxc8fcuxbcc0-scalar-uxc804uxd30cuxc640-psf}}

출력 좌표 \((x,y)\)는 \(\mathbf q_0\)를 중심으로 chief에 수직인 가상
평면의 mm 좌표다. 파장 \(\lambda_m=\lambda_0/n_i\)에서 국소 회절장은
전체 상수 위상을 제외하면

\[
U(x,y)\simeq\frac{1}{i\lambda_m}
 \iint P(\chi,\eta)
 \exp\left[-i\frac{2\pi}{\lambda_m}(x\chi+y\eta)\right]
 \,d\chi\,d\eta.
\tag{22}
\]

입력 적분의 uniform grid를 FFT로 계산하고, 출력의 intensity PSF를 면적
적분 1로 정규화한다.

\[
h(x,y)=\frac{|U(x,y)|^2}{\iint|U|^2\,dx\,dy},\qquad
\iint h\,dx\,dy=1.
\tag{23}
\]

주파수 축은 DFT 축을 실제 길이로 바꾼다. 예를 들어 padded 표본 수
\(N_p\)와 pupil 간격 \(\Delta\chi\)이면

\[
\Delta x=\frac{\lambda_m}{N_p\Delta\chi},\qquad
X_{\rm period}=\frac{\lambda_m}{\Delta\chi}.
\tag{24}
\]

따라서 zero padding은 PSF 표본 간격을 줄이지만 원래 파면의 정보를
추가하지 않는다. Ray 표본, 파면 grid, padding은 서로 독립적인 해상도
설정이다.

Zemax FFT PSF도 chief-normal 계산면을 명시한다. 따라서 위 \((x,y)\)는
일반적으로 global-\(z\) 센서의 spot 좌표와 같은 면이 아니다. 현재 엔진의
비축 센서 변환은 명시적으로 선택하는 국소 투영 근사이며, FFT 지표와 실제
센서면 영상은 구분한다. {[}R9{]}

\hypertarget{otfuxc640-mtf}{%
\subsubsection{5.3 OTF와 MTF}\label{otfuxc640-mtf}}

비간섭 결상의 complex OTF는 intensity PSF의 Fourier 변환이며 MTF는 그
절댓값이다. {[}R8{]}

\[
H(f_x,f_y)=\iint h(x,y)e^{-i2\pi(f_xx+f_yy)}\,dx\,dy,
\qquad \mathrm{MTF}=|H|,\qquad H(0,0)=1.
\tag{25}
\]

Focal 주파수는 image-space lp/mm로 보고한다. 요청 구간 0--120 lp/mm는
설정값이며 물리 cutoff를 뜻하지 않는다. 본 구현의 FFT S 곡선은 \(f_x\)
절단, T 곡선은 \(f_y\) 절단이다. 이 이름은 3절의 field-meridional 기하
T/S와 항상 일치하지 않는다. {[}R12{]}

각 field에서 T/S와 동일 pupil의 diffraction reference를 각각 반환한다.
주파수 방향과 줄무늬가 뻗는 방향은 수직이다. Zemax FFT의 T는 object-X
방향 줄무늬, S는 object-Y 방향 줄무늬에 대한 고정 이름이며 field 방위각에
따라 자동으로 radial 축을 선택하는 규약이 아니다. 본 엔진은 4.1절의
chief-normal 직교 frame을 사용하고, 실제 방향은
\texttt{mtf\_frequency\_axes\_global} 및
\texttt{mtf\_target\_line\_axes\_global}에 보존한다. 임의 2D field에서
제품의 pupil mapping까지 동일하다는 뜻은 아니다. {[}R12{]}

주파수 보간과 파장 합성은 complex \(H\)에 적용하고 마지막에 절댓값을
취한다. PSF의 위치 변화는 OTF 위상에 들어가므로, 파장별 MTF 절댓값을
단순 평균하면 횡색수차의 영향을 잃는다.

\hypertarget{uxac19uxc740-pupiluxc758-diffraction-reference}{%
\subsubsection{5.4 같은 pupil의 diffraction
reference}\label{uxac19uxc740-pupiluxc758-diffraction-reference}}

실제 계산과 \textbf{동일한 support·진폭·sampling}에서 \(W=0\)으로 놓아
\(P_{\rm DL}=|P|\)를 만든 뒤 같은 FFT를 반복한다. 이는 수차 위상만
제거한 비교 기준이다. 모든 field에 하나의 원형 aperture 공식을 적용하는
방식과 구별한다.

균일한 무차폐 원형 Fourier pupil은 독립 해석 기준을 갖는다. Radius가
\(a_\chi\)이면 \(f_c=2a_\chi/\lambda_m=2\mathrm{NA}/\lambda_0\)이고

\[
\mathrm{MTF}_{\rm circ}(f)=
\begin{cases}
\displaystyle\frac{2}{\pi}\left[\cos^{-1}\nu-\nu\sqrt{1-\nu^2}\right],&0\le\nu<1,\\
0,&\nu\ge1,
\end{cases}
\qquad\nu=\frac{|f|}{f_c}.
\tag{26}
\]

이 식은 원형 pupil의 정규화 자기상관으로 직접 얻는다. 실제 pupil이
다르면 같은-pupil FFT reference를 사용한다. 제품의 reference-field
선택과 이 엔진의 field별 비교 규약은 별도다. {[}R8, R12{]}

\hypertarget{uxbb34uxd55cuxb300-uxcd9cuxb825-uxac01uxb3c4-uxbd84uxd3ecuxc640-angular-mtf}{%
\subsubsection{5.5 무한대 출력: 각도 분포와 angular
MTF}\label{uxbb34uxd55cuxb300-uxcd9cuxb825-uxac01uxb3c4-uxbd84uxd3ecuxc640-angular-mtf}}

Afocal에서는 경계면의 실제 \((u,v)\)와
\(a_\Pi\exp[i2\pi(S_\Pi-S_c)/\lambda_0]\)를 Fourier 변환한다. 변환
주파수 \((\nu_u,\nu_v)\)가 주는 출력 방향 코사인은 \(l=\lambda_m\nu_u\),
\(m=\lambda_m\nu_v\)다. \(\gamma=\sqrt{1-l^2-m^2}\)일 때 {[}R7{]}

\[
F(l,m)=\iint U_\Pi(u,v)
 e^{-i2\pi(lu+mv)/\lambda_m}\,du\,dv,\qquad
d\Omega=\frac{dl\,dm}{\gamma}.
\tag{27}
\]

Scalar far-field obliquity를 포함한 정규화는 다음과 같다.
\(\mathcal D\)는 계산에서 표현한 propagating 방향 영역이다.

\[
\mathcal N=\iint_{\mathcal D}\gamma|F|^2\,dl\,dm,\qquad
h_\Omega(l,m)=\frac{\gamma^2|F|^2}{\mathcal N},\qquad
\iint_{\mathcal D}h_\Omega\,d\Omega=1.
\tag{28}
\]

Angular MTF는 \(l,m\)의 DFT 주파수를 이름만 바꾼 값이 아니다. 실제 각도
\(\alpha=\operatorname{atan2}(l,\gamma)\),
\(\beta=\operatorname{atan2}(m,\gamma)\)를 사용한다.

\[
H_\alpha(f)=\int h_\Omega(\Omega)e^{-i2\pi f\alpha(\Omega)}\,d\Omega
 \simeq\sum_kp_k e^{-i2\pi f\alpha_k},\qquad
p_k\simeq\frac{\gamma_k|F_k|^2\Delta l\Delta m}{\mathcal N}.
\tag{29}
\]

\(f\)의 단위는 cycles/rad다. FFT 방향 표본과 요청 \(f\)가 만드는 적분
위상을 함께 해상해야 한다. Focal과 angular 주파수를 연결하려면 지정한
물체--영상 또는 물체--각도 mapping의 국소 Jacobian이 필요하다. 단위만
바꿔 같은 숫자를 사용할 수 없다.

\hypertarget{uxad11uxb7c9uxacfc-riuxb294-radiance-uxbc0f-uxc804uxb825-uxbcf4uxc874uxc73cuxb85c-uxacc4uxc0b0uxd55cuxb2e4}{%
\subsection{6. 광량과 RI는 radiance 및 전력 보존으로
계산한다}\label{uxad11uxb7c9uxacfc-riuxb294-radiance-uxbc0f-uxc804uxb825-uxbcf4uxc874uxc73cuxb85c-uxacc4uxc0b0uxd55cuxb2e4}}

\hypertarget{lambertian-uxad11uxc6d0uxacfc-uxd22cuxacfcuxc728}{%
\subsubsection{6.1 Lambertian 광원과
투과율}\label{lambertian-uxad11uxc6d0uxacfc-uxd22cuxacfcuxc728}}

Lambertian은 방출 방향에 대한 radiance \(L_s\)가 일정하다는 뜻이다. 면적
\(dA_s\)가 방향 \(d\Omega_s\)로 보내는 전력은

\[
d\Phi=L_s\cos\theta_s\,dA_s\,d\Omega_s,
\qquad\Phi_{\rm hemi}=\pi L_s A_s.
\tag{30}
\]

본 광원의 연속 근사는 12.8×3.2 mm, fill factor 0.64다. \(L_s\)가 발광
active-area radiance이면 면적 평균은 \(0.64L_s\)다. 이미 panel-average
radiance이면 fill을 다시 곱하지 않는다. 320×80 픽셀 정보를 가지고 모든
발광 픽셀을 개별 추적할 필요는 없다. {[}R13{]}

실제 전력 투과율은 경로의 계면 투과와 내부 흡수를 결합한다. 무코팅·실수
굴절률 계면의 Fresnel 전력 반사율은

\[
R_s=\left(\frac{n_1c_i-n_2c_t}{n_1c_i+n_2c_t}\right)^2,
\quad R_p=\left(\frac{n_2c_i-n_1c_t}{n_2c_i+n_1c_t}\right)^2,
\quad T_s=1-R_s,\quad T_p=1-R_p.
\tag{31}
\]

단일 계면의 비편광 전력은 \((T_s+T_p)/2\)다. 기본 다면 계산은 서로
비간섭인 두 입력 편광 상태를 각 면의 s/p basis로 수송해 전력을 평균한다.
면마다 비편광으로 재설정한 단순 투과율 곱은 별도 scalar 옵션이다. 이
전력 계산 자체가 vector diffraction이나 코팅 위상 모델인 것은 아니다.
{[}R2, R14{]}

내부 전력 감쇠는 \(\tau_{\rm bulk}=\exp(-\mu\ell)\)다. 기준 두께
\(\ell_0\)의 내부 투과율 \(T_{\rm int}\)가 주어지면
\(\mu=-\ln T_{\rm int}/\ell_0\)로 환산한다. 이 값에는 계면 Fresnel
반사를 중복 포함하지 않는다.

\hypertarget{uxc77cuxbc18-uxacb0uxc0c1-ri}{%
\subsubsection{6.2 일반 결상 RI}\label{uxc77cuxbc18-uxacb0uxc0c1-ri}}

균일 장면의 각 영상점에서 유효 입사 방향을 역추적한다. Radiance의
\(L/n^2\) 불변성과 전력 투과율 \(\tau\)를 이용하면 {[}R13, R14{]}

\[
E(\mathbf x)=\left(\frac{n_i}{n_s}\right)^2L_s
 \int_{\Omega_{\rm accepted}}\tau\cos\theta_i\,d\Omega_i.
\tag{32}
\]

평면 영상의 direction cosines \((l,m)\)에서는
\(\cos\theta_i\,d\Omega_i=dl\,dm\)이므로 구현은 이 면적 측도에서
적분한다. 추가 cosine 또는 \(\cos^4\) 보정을 또 곱하지 않는다.

\[
E(\mathbf x)=\left(\frac{n_i}{n_s}\right)^2L_s
 \iint_{\mathcal P(\mathbf x)}\tau(l,m)\,dl\,dm,
\qquad RI(\mathbf x)=\frac{E(\mathbf x)}{\max_{\mathbf x\in\mathcal F}E(\mathbf x)}.
\tag{33}
\]

\(\mathcal F\)는 평가한 field 집합이므로 구현의 분모는 sampled peak다.
축상 대비 값도 별도로 구할 수 있다. 이 정의는 균일 Lambertian 장면과
최대 조도 정규화를 사용하는 Zemax RI의 해석 목적을 따른다. {[}R15{]}

\hypertarget{uxbb34uxd55cuxb300-uxcd9cuxb825uxc758-uxc0c1uxb300-uxbd84uxd3ec}{%
\subsubsection{6.3 무한대 출력의 상대
분포}\label{uxbb34uxd55cuxb300-uxcd9cuxb825uxc758-uxc0c1uxb300-uxbd84uxd3ec}}

실제 유한 발광면에서는 source 위치·방향에 식 (30)의 전력을 부여하고,
real-ray 경로의 투과율을 곱해 출사 방향 bin에 누적한다. 각 bin의 광도는

\[
\Phi_b=\sum_{k\in b}\Phi_k\tau_k,\qquad
I_b=\frac{\Phi_b}{\Delta\Omega_b},\qquad
I_{{\rm rel},b}=\frac{I_b}{\max_b I_b}.
\tag{34}
\]

이는 W/sr 또는 상대 광도이며, 영상면의 W/mm² RI와 다르다. 방출 전력 대비
총 통과 전력은 별도의 효율이다. 절대 radiance가 없으면 상대량으로
저장한다. 순수 real-ray 방향 histogram에는 회절이 없으며, 회절 angular
kernel을 원하면 다음 절의 광원--경계장 연결을 사용한다.

\hypertarget{uxc2e4uxc81c-uxad11uxc6d0uxc758-uxd68cuxc808uxacfc-uxc601uxc0c1-uxd569uxc131}{%
\subsection{7. 실제 광원의 회절과 영상
합성}\label{uxc2e4uxc81c-uxad11uxc6d0uxc758-uxd68cuxc808uxacfc-uxc601uxc0c1-uxd569uxc131}}

\hypertarget{uxc218uxc9d1-uxc804uxb825uxc5d0uxc11c-uxacbduxacc4uxc7a5-uxc9c4uxd3eduxc73cuxb85c}{%
\subsubsection{7.1 수집 전력에서 경계장
진폭으로}\label{uxc218uxc9d1-uxc804uxb825uxc5d0uxc11c-uxacbduxacc4uxc7a5-uxc9c4uxd3eduxc73cuxb85c}}

Source 점 \(\mathbf q\)의 단위 면적당 수집 전력을 \(Q(\mathbf q)\)라
한다. 이것은 입사 방향에 대한 적분이다.

\[
Q(\mathbf q)=L_{\rm panel}(\mathbf q)
 \int_{\Omega_{\rm accepted}(\mathbf q)}\tau(\mathbf q,\Omega_s)
 \cos\theta_s\,d\Omega_s.
\tag{35}
\]

실제 구현은 STOP 표본의 삼각형들을 source 방향 구면과 출력 경계면에 각각
mapping한다. Source에서의 projected solid angle을
\(\Delta\Omega_{p,j}=\int_j\cos\theta_s d\Omega_s\), 삼각형의 평균
투과율을 \(\overline\tau_j\)라 하면, source macro-cell 면적
\(\Delta A_s\)가 경계 삼각형으로 보내는 전력은

\[
P_j\simeq L_{\rm panel}\,\Delta A_s\,
 \overline\tau_j\,\Delta\Omega_{p,j},\qquad
Q\simeq\frac{\sum_jP_j}{\Delta A_s}.
\tag{36}
\]

경계 삼각형 면적을 \(\Delta A_{\Pi,j}\), 실제 ray와 경계 법선의 평균
cosine을 \(\overline c_j\)라 하면 전기장 진폭은 flux 관계로 정한다. SI
면적을 사용할 때

\[
P_j\simeq\frac{n_i}{2Z_0}|a_{\Pi,j}|^2\overline c_j\Delta A_{\Pi,j},
\qquad
|a_{\Pi,j}|\simeq
 \sqrt{\frac{2Z_0P_j}{n_i\overline c_j\Delta A_{\Pi,j}}}.
\tag{37}
\]

\(Z_0\)는 진공 임피던스다. 면적이 mm²이면 \(10^{-6}\)을 곱해 m²로
바꾼다. 이렇게 얻은 \(a_\Pi\)를 식 (21)에 연결한다. Source projected
solid angle, 경계면 면적, 실제 ray cosine을 모두 거치므로
\(\sqrt{P_j}\)를 그대로 pupil에 넣는 방식과 다르다. 이 삼각형 유한체적
근사는 본 엔진의 수송--회절 연결 규약이다.

PSF를 정규화하기 전의 ray 수집 전력, 경계장 flux, propagating wave
power를 따로 보존한다. 정규화 후 에너지 합이 맞는다는 사실만으로 경계
진폭의 정확성이 검증되지는 않는다.

\hypertarget{uxbe44uxac04uxc12d-uxd569uxc131uxacfc-uxacf5uxd1b5-uxc601uxc0c1-uxc88cuxd45c}{%
\subsubsection{7.2 비간섭 합성과 공통 영상
좌표}\label{uxbe44uxac04uxc12d-uxd569uxc131uxacfc-uxacf5uxd1b5-uxc601uxc0c1-uxc88cuxd45c}}

한 물점·한 파장은 pupil 내에서 복소장으로 전파한다. 서로 다른 LED 물점과
서로 다른 파장은 intensity를 합한다. Source별 정규화 intensity kernel을
\(h(\mathbf x;\mathbf q)\)라 하면

\[
E_{\rm image}(\mathbf x)=\int Q(\mathbf q)h(\mathbf x;\mathbf q)
 B_{\rm image}(\mathbf q)\,dA_s.
\tag{38}
\]

\(B_{\rm image}\)는 샘플 이미지의 \textbf{선형 방출량 modulation}이다.
sRGB 입력은 먼저 선형값으로 변환한다. Source cell별 \(Q\)에는 fill, 방향
수집률과 투과가 이미 들어 있으므로 합성 단계에서 다시 적용하지 않는다.

출력 pixel \(b\)의 kernel 확률질량
\(h_{jb}=\int_bh(\mathbf x;\mathbf q_j)dA_i\)로 이산화하면

\[
\Phi_b=\sum_jQ_j\Delta A_{s,j}B_{{\rm image},j}h_{jb},
\qquad E_b=\frac{\Phi_b}{\Delta A_{i,b}}.
\tag{39}
\]

모든 kernel은 같은 물리 detector grid에 등록한다. Crop 밖으로 나간 PSF
tail은 누락 전력으로 보존하며 저장된 부분만 재정규화하지 않는다. 공간에
따라 mapping과 PSF가 달라지므로 일반적으로 전체 이미지를 한 PSF로
convolution하는 것으로 충분하지 않다.

RGB도 같은 global grid에서 합한다. Local OTF 비교가 필요하면 공통
주파수와 registration 위상을 유지한 상태로

\[
H_{\rm poly}=\frac{\sum_\lambda w_\lambda H_\lambda}{\sum_\lambda w_\lambda},
\qquad\mathrm{MTF}_{\rm poly}=|H_{\rm poly}|.
\tag{40}
\]

G 단색 대표 MTF에는 이 합성이 필요 없다. 광원 pixel pitch나 fill의
전달함수도 lens-only MTF에 자동 결합하지 않는다.

\hypertarget{uxade0uxc77c-uxc628uxb3c4-uxbcc0uxd654uxb294-uxc0c8-uxcc98uxbc29uxc73cuxb85c-uxac19uxc740-uxc5d4uxc9c4uxc744-uxbc18uxbcf5uxd55cuxb2e4}{%
\subsection{8. 균일 온도 변화는 새 처방으로 같은 엔진을
반복한다}\label{uxade0uxc77c-uxc628uxb3c4-uxbcc0uxd654uxb294-uxc0c8-uxcc98uxbc29uxc73cuxb85c-uxac19uxc740-uxc5d4uxc9c4uxc744-uxbc18uxbcf5uxd55cuxb2e4}}

온도 해석은 \textbf{굴절률 변화와 자유 열팽창으로 온도별 광학 상태를
만든 뒤 동일한 추적·성능 계산을 반복}하는 것이다. 렌즈 내부 공간 온도
분포, 응력, GRIN, 복굴절, mount 구조 해석은 이 모델의 입력이 아니다.

\hypertarget{uxad74uxc808uxb960uxacfc-uxd615uxc0c1uxc758-uxc628uxb3c4-uxc758uxc874uxc131}{%
\subsubsection{8.1 굴절률과 형상의 온도
의존성}\label{uxad74uxc808uxb960uxacfc-uxd615uxc0c1uxc758-uxc628uxb3c4-uxc758uxc874uxc131}}

공급된 \(n(\lambda_0,T)\) 표를 평가하거나, 제공된 absolute
\(\beta(\lambda_0)=dn/dT\)의 선형 모델을 지정하면

\[
n(\lambda_0,T)=n(\lambda_0,T_c)+\beta(\lambda_0)(T-T_c).
\tag{41}
\]

\(T_c\)는 굴절률 자료의 기준 온도다. 선형 모델의 유효범위는 자료의
범위를 따른다. SCHOTT TIE-19 등 제조사 모델은 absolute/relative
thermo-optic 계수 구분의 근거이며, 현재 엔진이 모든 제조사 catalog 식을
자동 지원한다는 뜻은 아니다. 필요한 \(n(T)\) 데이터로 변환해 제공할 수
있다. {[}R11{]}

형상 기준온도 \(T_0\)에 대한 등방 배율을 \(s(T;T_0)\)라 한다. 기준온도와
CTE 기준이 같은 선형 모델은 \(s=1+\alpha(T-T_0)\)다. Instantaneous CTE
\(\alpha=(1/L)dL/dT\)를 제공하면

\[
s(T;T_0)=\exp\left(\int_{T_0}^{T}\alpha(\vartheta)\,d\vartheta\right).
\tag{42}
\]

CTE의 정의를 선택한 뒤 다음 형상 항등식을 적용한다.

\[
z_{s,T}(r)=s\,z_{s,0}(r/s),\qquad
R_T=sR_0,\quad c_T=c_0/s,\quad k_T=k_0,
\quad A_{2m,T}=A_{2m,0}s^{1-2m}.
\tag{43}
\]

중심 두께와 clear/mechanical diameter도 \(s\)배다. 식 (43)은 고차
asphere를 포함한 전체 표면의 기하 similarity에서 직접 나온다. 렌즈별
anchor와 air gap 정책으로 각 vertex 위치를 정한다. STOP 반경 유지/렌즈
추종, OBJ/IMA 위치 유지 여부도 명시한다. 모든 온도 상태는 원 기준
처방에서 독립 생성한다.

\hypertarget{uxace0uxc815uxba74uxacfc-uxacf5uxd1b5-uxc7acuxcd08uxc810}{%
\subsubsection{8.2 고정면과 공통
재초점}\label{uxace0uxc815uxba74uxacfc-uxacf5uxd1b5-uxc7acuxcd08uxc810}}

기본 비교는 지정한 IMA를 유지하여 focus drift, spot, MTF를 계산한다.
재초점을 허용하면 \textbf{모든 field에 공통인 하나의 영상면}을 이동한다.

최종 균질 매질에서 ray line은
\(\mathbf x_{hi}(\delta)=\mathbf a_{hi}+\delta\mathbf b_{hi}\)다. 각
field \(h\)의 pupil 가중치를 정규화하고
\(\Delta\mathbf a,\Delta\mathbf b\)를 해당 field의 평균에서 뺀 값으로
정의하면, field weight \(v_h\)에 대한 geometric centroid-RMS² 최소점은

\[
\delta_*=-\frac{\sum_hv_h\sum_i\widehat w_{hi}
 \Delta\mathbf a_{hi}\cdot\Delta\mathbf b_{hi}}
 {\sum_hv_h\sum_i\widehat w_{hi}\|\Delta\mathbf b_{hi}\|^2}.
\tag{44}
\]

이 식은 고정된 ray 집합과 가중치에 대한 quadratic merit의 직접 미분
결과다. MTF 최대 초점과 같은 목적함수는 아니다. MTF focus sweep은 선택한
공통 IMA 위치마다 식 (19) 이후를 다시 계산해 비교한다.

IMA만 이동할 때는 같은 경계면 \(\Pi\)와 그 위의 물리 진폭 \(a_\Pi\)를
유지한다. Lens-only 최초 기준거리 \(L_0\)에서 \(B=1\)이었다면
\(a_\Pi=L_0/(L_0^2+u^2+v^2)\)이며, 이동 후에는 새 \(L\)과 식 (21)로
\(B\)를 다시 구성한다. 매번 \(B=1\)로 초기화하면 focus뿐 아니라 조명도
바뀐다. 처방·온도·STOP·물점·파장이 바뀌면 그 상태의 광선을 새로
추적한다.

\hypertarget{uxc774-uxbb38uxc11cuxc758-uxc6d0uxb9acuxb97c-uxcf54uxb4dcuxc5d0uxc11c-uxcc3euxb294-uxc704uxce58}{%
\subsection{9. 이 문서의 원리를 코드에서 찾는
위치}\label{uxc774-uxbb38uxc11cuxc758-uxc6d0uxb9acuxb97c-uxcf54uxb4dcuxc5d0uxc11c-uxcc3euxb294-uxc704uxce58}}

\begin{longtable}[]{@{}
  >{\raggedright\arraybackslash}p{(\columnwidth - 2\tabcolsep) * \real{0.5000}}
  >{\raggedright\arraybackslash}p{(\columnwidth - 2\tabcolsep) * \real{0.5000}}@{}}
\toprule\noalign{}
\begin{minipage}[b]{\linewidth}\raggedright
계산 단계
\end{minipage} & \begin{minipage}[b]{\linewidth}\raggedright
구현 모듈·함수
\end{minipage} \\
\midrule\noalign{}
\endhead
\bottomrule\noalign{}
\endlastfoot
입력·설정·선택적 물성 & \texttt{models.py}, \texttt{settings.py},
\texttt{materials.py}, \texttt{io.py} \\
sag, normal, 교점, Snell, STOP aiming & \texttt{geometry.py}:
\texttt{sag}, \texttt{gradient}, \texttt{intersect}, \texttt{refract},
\texttt{aim\_ray}, \texttt{trace\_bundle} \\
EFL/BFL와 aperture 변환 & \texttt{paraxial.py}:
\texttt{paraxial\_properties}, \texttt{resolve\_aperture} \\
spot·mapping·수차·공통 초점 & \texttt{metrics.py} \\
경계면·OPD·FFT·OTF·회절 기준 & \texttt{wave.py}:
\texttt{make\_boundary\_state}, \texttt{analyze\_wave},
\texttt{fft\_psf}, \texttt{angular\_otf} \\
광량·RI·원거리 수송 & \texttt{radiometry.py} \\
물리 source에서 회절장으로 연결 & \texttt{source\_wave.py}:
\texttt{source\_diffraction\_kernel} \\
source별 kernel의 비간섭 합성 & \texttt{imaging.py}:
\texttt{synthesize\_image}, \texttt{registered\_spectral\_sum} \\
균일 열 상태 생성 & \texttt{thermal.py}: \texttt{scale\_surface},
\texttt{thermal\_system} \\
선택한 분석의 실행 & \texttt{api.py}: \texttt{evaluate} \\
\end{longtable}

현재 모델은 동축 회전대칭 면·균질 등방 매질·순차 직접 투과·scalar 회절을
대상으로 한다. 수치 수렴과 이 물리 모델의 적용성은 각각 확인한다. 상세
알고리즘, 예외 처리, 기존 논쟁 및 코드 검증 이력은
\texttt{calculation\_spec\_ko.md}, 루트 검증 문서와 역할별 리뷰를
참고한다. 본 문서의 식과 테스트 통과만으로 상용툴 상관 검증이나 실제
측정 정확도가 입증되지는 않는다.

\hypertarget{uxcc38uxace0uxbb38uxd5ccuxacfc-uxc0acuxc6a9uxd55c-uxbc94uxc704}{%
\subsection{10. 참고문헌과 사용한
범위}\label{uxcc38uxace0uxbb38uxd5ccuxacfc-uxc0acuxc6a9uxd55c-uxbc94uxc704}}

모든 웹 원문 확인일: 2026-10-18. 아래 자료는 원저자·대학 강의·제조사
공식 문서다. 문서의 자체 좌표변환, 정규화, 구현 계약과 직접 유도는 각
절에서 구분했다. 제품 문서는 기능 정의의 근거이며 이 코드의 구현 검증을
대신하지 않는다.

\begin{itemize}
\tightlist
\item
  \textbf{{[}R1{]} Ansys, Even Asphere.} Sag와 even-power 계수 정의.
  \href{https://ansyshelp.ansys.com/public/Views/Secured/Zemax/v261/en/OpticStudio_User_Guide/OpticStudio_Help/topics/Even_Asphere.html}{공식
  문서}
\item
  \textbf{{[}R2{]} Pharr, Jakob, Humphreys, PBRT 4e, Specular Reflection
  and Transmission.} Snell·Fresnel 및 전력/진폭 구분.
  \href{https://pbr-book.org/4ed/Reflection_Models/Specular_Reflection_and_Transmission}{원저자
  웹북}
\item
  \textbf{{[}R3{]} Ansys / Nam-Hyong Kim, How to use Ray Aiming.} 실제
  STOP 목표를 맞추는 반복 추적과 pupil의 의미.
  \href{https://optics.ansys.com/hc/en-us/articles/42661778056083-How-to-use-Ray-Aiming}{공식
  기술 문서}
\item
  \textbf{{[}R4{]} Michael J. Hayford, ray-optics, parax.firstorder.}
  Reduced-coordinate first-order 광학의 공개 구현. 식 (12)--(13)은 본
  문서의 좌표 정의에서 직접 도출한다.
  \href{https://ray-optics.readthedocs.io/en/stable/_modules/rayoptics/parax/firstorder.html}{원저자
  소스 문서}
\item
  \textbf{{[}R5{]} Ansys, Field Curvature and Distortion, Discussion.}
  Parabasal T/S focus와 실제 small-field chief 기준 distortion.
  \href{https://ansyshelp.ansys.com/public/Views/Secured/Zemax/v252/en/OpticStudio_User_Guide/OpticStudio_Help/topics/Field_Curvature_and_Distortion.html}{공식
  문서}
\item
  \textbf{{[}R6{]} TU Delft, Advanced Optical Imaging, §8.3 Zernike
  polynomials, 일반식 (17)--(19).} 원판 basis와 정규화. 단일-index
  예시표 대신 본 코드의 \((n,m,0/c/s)\) tuple 정의를 사용한다.
  \href{https://qiweb.tudelft.nl/aoi/wavefieldaberrations/wavefieldaberrations/}{대학
  강의 원자료}
\item
  \textbf{{[}R7{]} Robert L. Lucke, Rayleigh-Sommerfeld Fraunhofer
  Diffraction, 2006, §§2--3, 식 (5), (7), (8).} Obliquity를 포함한
  Cartesian--direction-cosine Fourier 관계. 식 (19)--(21)의 실제 ray
  boundary 연결은 본 구현의 유도다.
  \href{https://arxiv.org/pdf/physics/0604229}{논문}
\item
  \textbf{{[}R8{]} MIT OCW, 2.71/2.710 Optics, Lecture 22, 2009, 강의
  슬라이드 11--20.} 비간섭 intensity PSF, OTF, pupil 자기상관과 MTF.
  \href{https://ocw.mit.edu/courses/2-71-optics-spring-2009/bebe78c6a2ae2603a0d1a3ea85e0b4ca_MIT2_71S09_lec22.pdf}{강의
  PDF}
\item
  \textbf{{[}R9{]} Ansys, FFT PSF, Assumptions / Discussion.}
  Chief-normal 가상면, scalar/Fraunhofer 범위와 sampling 의미.
  \href{https://ansyshelp.ansys.com/public/Views/Secured/Zemax/v251/en/OpticStudio_User_Guide/OpticStudio_Help/topics/FFT_PSF.html}{공식
  문서}
\item
  \textbf{{[}R10{]} SCHOTT, TIE-29: Refractive Index and Dispersion,
  §§2.1--2.3.} Abbe 수와 Sellmeier 분산.
  \href{https://media.schott.com/api/public/content/aaa572afd854434fb7b3faa4bc46103f?v=c0f4fa52}{제조사
  PDF}
\item
  \textbf{{[}R11{]} SCHOTT, TIE-19: Temperature Coefficient of the
  Refractive Index, §§1--2.} Absolute/relative index, thermo-optic
  계수와 기준 상태.
  \href{https://media.schott.com/api/public/content/63b62a0d73a640efa8c6fde92f2ad5c2?v=97a6a05a}{제조사
  PDF}
\item
  \textbf{{[}R12{]} Ansys, FFT MTF, Discussion.} Image-space 주파수, FFT
  T/S 규약, diffraction reference와 pupil sampling.
  \href{https://ansyshelp.ansys.com/public/Views/Secured/Zemax/v252/en/OpticStudio_User_Guide/OpticStudio_Help/topics/FFT_MTF.html}{공식
  문서}
\item
  \textbf{{[}R13{]} Pharr, Jakob, Humphreys, PBRT 3e, §5.5 Working with
  Radiometric Integrals.} Radiance 적분, projected solid angle,
  면적--입체각 변환.
  \href{https://www.pbr-book.org/3ed-2018/Color_and_Radiometry/Working_with_Radiometric_Integrals}{원저자
  웹북}
\item
  \textbf{{[}R14{]} Pharr, Jakob, Humphreys, PBRT 4e, Dielectric BSDF.}
  매질 경계의 radiance 변환과 power 수송.
  \href{https://pbr-book.org/4ed/Reflection_Models/Dielectric_BSDF}{원저자
  웹북}
\item
  \textbf{{[}R15{]} Ansys, Relative Illumination, Discussion.} 균일
  Lambertian 장면, 최대 조도 정규화, image-direction-cosine 적분.
  \href{https://ansyshelp.ansys.com/public/Views/Secured/Zemax/v26102/en/OpticStudio_User_Guide/OpticStudio_Help/topics/Relative_Illumination.html}{공식
  문서}
\end{itemize}

\end{document}
